%% notes-tex/010-additive-baselines/sections/8-modernize.tex

\section{Modernizing the 010 build\secstatus{todo}}
\label{sec:modernize}

Two of the obstacles in this document have the same root. Per-record predictions
had to be rebuilt by hand rather than recomputed through the stock evaluation
path, and the transformer has no row in Table~\ref{tab:disjoint}, because the
December 2025 build no longer loads under the current code.

\subsection{What actually broke\secstatus{todo}}

Two later changes invalidated it, and neither is about what was measured.

\begin{enumerate}
\item The perturbation-ontology refactor of 2026-07-08 renamed the deletion
  literals. Of the 1{,}130{,}196 stored perturbations, 1{,}008{,}761 carry
  \texttt{perturbation\_type: "deletion"} with \texttt{deletion\_type: "KanMX"},
  which is now \texttt{sga\_kanmx\_deletion}, and 38{,}742 carry
  \texttt{deletion\_type: "mean"}, which is now \texttt{mean\_deletion}. The
  remaining 82{,}693 are allele and temperature-sensitive-allele perturbations
  and were not renamed. No record in this build is a \texttt{natMX} deletion.
\item The component-based media change of 2026-07-14 made
  \texttt{Media.is\_synthetic} required. Every record and every reference in the
  build lacks it.
\end{enumerate}

An earlier version of this document named only the first, and named the wrong
replacement vocabulary for it. Both are pure vocabulary changes: neither alters
which gene was perturbed, what was measured, or what the label is.

\subsection{The migration, and why it copies rather than rebuilds\secstatus{todo}}

Every published 010 artifact addresses records by position. The split file, the
label table, the perturbed-gene index and the saved per-record predictions are
all position-keyed, and the transformer checkpoints were selected against that
split. A rebuild from the served graph returns records in a different order, so
it would preserve the science and destroy the correspondence. Precedent is on
record: carrying 010's splits into experiment 025 by a key that included the
perturbation type matched 682 records of 376{,}732, because the type strings had
been renamed in between, while a gene-set key matched all of them.

The migration is therefore a key-preserving copy. It rewrites the two vocabulary
fields, writes a sibling build, and leaves the frozen original untouched, which
it must, since the original's files are owned by the graph-build user and
hardlinked into the graph tree. All 376{,}732 records migrate, all 376{,}732
validate against the current schema, and the check that the copy is faithful is
that every position still carries the same three systematic gene names and the
same label.

One thing the migration deliberately does not fix. The frozen build records YEPD
at 30 degrees. The Kuzmin loaders now emit triple-mutant selection medium at 26
degrees, which is the corrected reading of the screen. Applying that here would
change the experimental content of the very records whose published results the
migration exists to preserve. A build with the corrected environment is a
different data version, and it is a rebuild.

\subsection{Which experiment ran on which data\secstatus{todo}}
\label{sec:versioning}

The migration creates a second 010 build, so the question of recording which
version an experiment used is now concrete rather than hypothetical.

A commit hash does not answer it, for a reason this document supplies. The
off-by-28 defect of Section~\ref{sec:offby28} happened with no code change at
all: the library was untouched between the runs, and what moved was the genome
gene set the cell graph is rebuilt from. Two runs at the same commit scored
different genes. Identity has to be carried by the built artifact, not by the
code that built it.

What already exists is most of the machinery. Dataset builds get a
\texttt{build\_manifest.json} recording a contract fingerprint for every schema
symbol the loader depends on, plus build time, host and commit, and staleness is
judged by refingerprinting rather than by comparing commits. The gap is that
experiment builds are not produced by that class and get no manifest, so neither
010 nor any later experiment build has one.

Three things would close it, in order of how much they are worth.

\begin{enumerate}
\item \textbf{Pin the index space with the build.} The gene ordering a
  checkpoint was trained against is the artifact whose drift caused the worst
  defect here, and it is currently recomputed at every use. Writing it into the
  build and loading it, rather than reconstructing it, makes the failure
  impossible instead of merely detected. The guard added in
  Section~\ref{sec:offby28} is the detection half.
\item \textbf{Give experiment builds a manifest.} Extend the existing manifest to
  the experiment-build class, so a build records the schema contract it was made
  against and any run can name its data version.
\item \textbf{Record the source graph's identity.} A build made by querying the
  graph database stores the query but nothing about which database answered it.
  A build identifier from the server, recorded alongside the query, would make
  the chain complete.
\end{enumerate}

The migrated build carries a \texttt{build\_provenance.json} naming its source,
that source's LMDB checksum, the exact field rewrites, and what was deliberately
not changed. That is the first item's shape, written by hand for one build.
